% ============ 第一章 §1 代数数域 (原书 pp.3-17) ============
\section{代数数域}\label{sec:ch1-s1}

有理数域 $\mathbb{Q}$ 的有限(次)扩域 $K$ 叫作是\emph{代数数域},简称作数域,这是代数数论的基本研究对象.如果扩张次数 $[K:\mathbb{Q}]$ 是 $n$,则 $K$ 也叫作是 $n$ 次(数)域.由于有限扩张必然是代数扩张,所以数域 $K$ 中每个元素均是 $\mathbb{Q}$ 上的代数元素.根据代数基本定理(附录 A, (18)),复数域 $\mathbb{C}$ 是 $\mathbb{Q}$ 的代数封闭扩域.从而数域 $K$ 中每个元素均可看成是复数,而每个数域 $K$ 均可看成是 $\mathbb{C}$ 的子域.如果 $K\subseteq\mathbb{R}$($\mathbb{R}$ 表示实数域),则称 $K$ 为实域,否则称 $K$ 为虚域.元素 $\alpha\in\mathbb{C}$ 如果是 $\mathbb{Q}$ 上的代数元素(即存在 $f(x)\in\mathbb{Q}[x]$,$\deg f(x)\geqslant 1$,使得 $f(\alpha)=0$),我们称 $\alpha$ 为代数数,否则便叫作是超越数.所有代数数全体构成域 $\Omega$,叫作是 $\mathbb{Q}$ 的代数闭包.事实上每个数域均是 $\Omega$ 的子域.而 $\mathbb{C}$ 是大于 $\Omega$ 的.换句话说,超越数是存在的.例如,可以证明 $\pi$ 和 $e$ 均是超越数,并且超越数比代数数还要多(习题 1).

关于域的代数扩张的一般事实请参见附录 A, III.在这一节里,我们就数域的情形再作一些补充.

\subsection{单扩张定理}\label{sec:ch1-s1-1}

设 $L/K$ 是数域的扩张(即 $L$ 和 $K$ 均是数域并且 $K\subseteq L$).由于扩张 $L/\mathbb{Q}$ 和 $K/\mathbb{Q}$ 均是有限的,从而 $L/K$ 也是有限扩张.令扩张次数为 $[L:K]=n$,而 $\omega_1,\ldots,\omega_n$ 是向量空间 $L$ 的一组 $K$-基,则 $L$ 中每个元素均可惟一地写为
\[
  k_1\omega_1+\cdots+k_n\omega_n,\quad k_i\in K.
\]
特别地有 $L=K(\omega_1,\omega_2,\cdots,\omega_n)$,即 $L/K$ 是有限生成扩张.我们现在要进一步证明:

\begin{theorem}\label{thm:1.1}
每个数域扩张 $L/K$ 均是单扩张,即存在 $\gamma\in L$,使得 $L=K(\gamma)$.
\end{theorem}

\begin{proof}
我们只要对 $L=K(\alpha,\beta)$ 的情形证明定理即可,因为一般情形 $L=K(\omega_1,\cdots,\omega_n)$ 可由此对 $n$ 归纳证得.现设 $L=K(\alpha,\beta)$.令 $f(x),g(x)\in K[x]$ 分别是元素 $\alpha$ 和 $\beta$ 在 $K$ 上的极小多项式,它们在 $\mathbb{C}[x]$ 中分解为
\[
  f(x)=\prod_{i=1}^{n}(x-\alpha_i),\quad
  g(x)=\prod_{j=1}^{m}(x-\beta_j),\quad \alpha_i,\beta_j\in\mathbb{C},
\]
其中 $n=\deg f$,$m=\deg g$.不妨设 $\alpha=\alpha_1$,$\beta=\beta_1$.由于 $f(x)$ 和 $g(x)$ 均是 $K[x]$ 中不可约多项式,从而它们均无重根,即 $\alpha_i$($1\leqslant i\leqslant n$)两两相异,而 $\beta_j$($1\leqslant j\leqslant m$)也两两相异.现在于有限集合
\[
  \left\{\frac{\alpha_i-\alpha}{\beta-\beta_j}\;\middle|\;1\leqslant i\leqslant n,\ 1\leqslant j\leqslant m,\ j\neq 1\right\}
\]
之外取一个非零有理数 $c$,不难看出 $n$ 个复数 $\alpha_i+c\beta$($1\leqslant i\leqslant n$)两两相异.令 $\gamma=\alpha+c\beta$,则多项式 $h(x)=f(\gamma-cx)$ 属于 $K(\gamma)[x]$,$h(\beta_1)=0$,而 $\beta_2,\cdots,\beta_m$ 均不为 $h(x)$ 的根.于是在 $\mathbb{C}[x]$ 中 $(h(x),g(x))=x-\beta_1$.注意域上两个多项式的最大公因子可以用辗转相除法求得,而这个过程在 $K(\gamma)[x]$ 中和 $\mathbb{C}[x]$ 中都是一样的,因此在 $K(\gamma)[x]$ 中也有 $(h(x),g(x))=x-\beta_1$.特别地 $x-\beta_1\in K(\gamma)[x]$,这就表明 $\beta=\beta_1\in K(\gamma)$,于是 $\alpha=\gamma-c\beta\in K(\gamma)$.从而 $K(\alpha,\beta)\subseteq K(\gamma)$.另一方面,由于 $\gamma=\alpha+c\beta$,$c\in K$,从而 $K(\gamma)\subseteq K(\alpha,\beta)$ 显然成立.这就证明了 $K(\alpha,\beta)=K(\gamma)$,从而也证明了定理 \ref{thm:1.1}.
\end{proof}

\subsection{数域的嵌入}\label{sec:ch1-s1-2}

设 $L/K$ 是数域的扩张,正如附录 A, III 中所述,每个域的单同态 $\sigma\colon L\to\mathbb{C}$ 均叫作 $L$ 在 $\mathbb{C}$ 中的一个嵌入.如果 $\sigma$ 在 $K$ 上的限制 $\sigma|_K$ 是域 $K$ 上的恒等自同构(即对每个 $k\in K$ 均有 $\sigma(k)=k$),则称 $\sigma$ 是 $K$-嵌入.利用上面的单扩张定理我们可以证明:$L$ 恰好有 $[L:K]$ 个 $K$-嵌入.事实上,我们可以证明下面更为一般的结论:

\begin{theorem}\label{thm:1.2}
设 $L/K$ 是数域的扩张,$[L:K]=n$,则每个嵌入 $\sigma\colon K\to\mathbb{C}$ 均可以 $n$ 种不同的方法扩充到 $L$ 上.换句话说,恰好存在 $n$ 个不同的嵌入 $\tau_i\colon L\to\mathbb{C}$($1\leqslant i\leqslant n$),使得 $\tau_i|_K=\sigma$.
\end{theorem}

\begin{proof}
由单扩张定理我们可以令 $L=K(\gamma)$.命
\[
  f(x)=c_0+c_1x+\cdots+c_{n-1}x^{n-1}+x^n\in K[x]
\]
是 $\gamma$ 在 $K$ 上的极小多项式,则 $\deg f=n$,而 $L$ 中元素均可惟一地表示成
\[
  \alpha=k_0+k_1\gamma+\cdots+k_{n-1}\gamma^{n-1},\quad k_i\in K
\]
(附录 A, (12) 及其注记).设 $\tau\colon L\to\mathbb{C}$ 是一个嵌入并且 $\tau|_K=\sigma$,则
\[
  \tau(\alpha)=\sigma(k_0)+\sigma(k_1)\tau(\gamma)+\cdots+\sigma(k_{n-1})\tau(\gamma)^{n-1},
\]
从而 $\tau$ 由它在 $\gamma$ 上的值所完全决定.考虑多项式
\[
  \sigma f(x)=\sigma(c_0)+\sigma(c_1)x+\cdots+\sigma(c_n)x^{n-1}+x^n\in\sigma(K)[x].
\]
由于 $\sigma\colon K\to\sigma(K)$ 是域的同构,不难看出 $\sigma f$ 是 $\sigma(K)[x]$ 中的 $n$ 次不可约多项式,从而它有 $n$ 个不同的复根 $\rho_1,\cdots,\rho_n$.由于
\[
  \sigma f(\tau(\gamma))=\sigma(c_0)+\sigma(c_1)\tau(\gamma)+\cdots+\sigma(c_{n-1})\tau(\gamma)^{n-1}+\tau(\gamma)^n=\tau(f(\gamma))=0,
\]
这就表明 $\tau(\gamma)$ 必为某个 $\rho_i$.从而 $\sigma$ 到 $L$ 上的扩充至多有 $n$ 个.现在对每个 $i$($1\leqslant i\leqslant n$),作映射
\[
  \tau_i\colon L\to\mathbb{C},\qquad
  \tau_i(k_0+k_1\gamma+\cdots+k_{n-1}\gamma^{n-1})
  =\sigma(k_0)+\sigma(k_1)\rho_i+\cdots+\sigma(k_{n-1})\rho_i^{n-1},
\]
容易验证这是域的同态.设 $k_0+k_1\gamma+\cdots+k_{n-1}\gamma^{n-1}\in\Ker\tau_i$(同态 $\tau_i$ 的核),则
\[
  \sigma(k_0)+\sigma(k_1)\rho_i+\cdots+\sigma(k_{n-1})\rho_i^{n-1}=0,
\]
从而
\[
  \sigma f(x)\mid\sigma(k_0)+\sigma(k_1)x+\cdots+\sigma(k_{n-1})x^{n-1},
\]
于是 $f(x)\mid k_0+k_1x+\cdots+k_{n-1}x^{n-1}$(为什么?).但是 $f(x)$ 为 $K[x]$ 中 $n$ 次不可约多项式,所以只能是 $k_0=k_1=\cdots=k_{n-1}=0$.这就表明 $\Ker\tau_i=(0)$,即 $\tau_i$ 是嵌入.又显然 $\tau_i|_K=\sigma$ 并且 $\tau_i(\gamma)=\rho_i$.而 $\rho_i$($1\leqslant i\leqslant n$)是两两相异的,从而 $\tau_i$($1\leqslant i\leqslant n$)是 $\sigma$ 到 $L$ 上的 $n$ 个不同的扩充.这就证明了定理 \ref{thm:1.2}.
\end{proof}

在定理 \ref{thm:1.2} 中特别取 $\sigma$ 为域 $K$ 的恒等自同构,我们就得到:

\begin{corollarynonum}
设 $L/K$ 为数域扩张,则从 $L$ 到 $\mathbb{C}$ 恰好有 $[L:K]$ 个不同的 $K$-嵌入.
\end{corollarynonum}

\begin{remark}
设 $L=K(\gamma)$,$f(x)\in K[x]$ 是 $\gamma$ 在 $K$ 上的极小多项式,$\deg f=n=[L:K]$,则 $f(x)$ 在 $\mathbb{C}$ 中有 $n$ 个不同的根 $\gamma_i$($1\leqslant i\leqslant n$),其中有一个根 $\gamma_1$ 为 $\gamma$.它们叫作是 $\gamma$ 的 $K$-共轭元素(附录 A, III).从定理 \ref{thm:1.2} 的证明可知映射
\[
  \tau_i\colon L=K(\gamma)\to K(\gamma_i),\quad \tau_i(\gamma)=\gamma_i\quad (1\leqslant i\leqslant n)
\]
就是 $L$ 到 $\mathbb{C}$ 中的全部 $n$ 个 $K$-嵌入.这是 $n$ 个不同的嵌入方式,因为元素 $\gamma$ 的象 $\gamma_i$($1\leqslant i\leqslant n$)彼此不同.对于 $i=1$,$K(\gamma_1)=L$ 并且 $\tau_1$ 是恒等自同构.$n$ 个与 $L$ 同构的域 $K(\gamma_i)$($1\leqslant i\leqslant n$)叫作 $L$ 的 $K$-共轭域,它们不必彼此不同.特别当这些域彼此相同,从而均为 $L=K(\gamma_1)$ 的时候,也就是 $\gamma_i\in L$($1\leqslant i\leqslant n$)的时候,$\tau_i$($1\leqslant i\leqslant n$)均是域 $L$ 的 $K$-自同构.这时称 $L/K$ 为伽罗瓦扩张,而 $\Gal(L/K)=\{\tau_1,\cdots,\tau_n\}$ 是 $L/K$ 的伽罗瓦群.而对于一般的情形,由于 $K(\gamma_1,\gamma_2,\cdots,\gamma_n)$ 是 $f(x)$ 在 $K$ 上的分裂域,从而 $K(\gamma_1,\cdots,\gamma_n)/K$ 是伽罗瓦扩张(附录 A, III, (15)).并且不难看出 $K(\gamma_1,\gamma_2,\cdots,\gamma_n)$ 是 $K$ 的包含 $L$ 的最小伽罗瓦扩张.称 $K(\gamma_1,\cdots,\gamma_n)$ 为扩张 $L/K$ 的\emph{正规闭包}.

如果 $K=\mathbb{Q}$,即 $L$ 是 $n=[L:\mathbb{Q}]$ 次数域,$L=\mathbb{Q}(\gamma)$.令 $f(x)\in\mathbb{Q}[x]$ 是 $\gamma$ 在 $\mathbb{Q}$ 上的极小多项式,则存在恰好 $n$ 个域的嵌入 $\tau_i\colon L=\mathbb{Q}(\gamma)\to\mathbb{Q}(\gamma_i)\subseteq\mathbb{C}$,使得 $\tau_i(\gamma)=\gamma_i$($1\leqslant i\leqslant n$),其中 $\gamma_i$($1\leqslant i\leqslant n$)是 $f(x)$ 的 $n$ 个不同的根(注意:数域的嵌入必为 $\mathbb{Q}$-嵌入!).不妨设前 $r_1$ 个是实根而后 $r_2$ 对是虚根,即
\[
  \gamma_i\in\mathbb{R}\quad (1\leqslant i\leqslant r_1),
\]
\[
  \gamma_{r_1+j}=\overline{\gamma}_{r_1+r_2+j}\notin\mathbb{R}\quad
  (1\leqslant j\leqslant r_2,\ r_1+2r_2=n).
\]
于是 $L$ 的前 $r_1$ 个共轭域 $\mathbb{Q}(\gamma_i)$ 为实域,我们称这 $r_1$ 个嵌入 $\tau_i\colon L=\mathbb{Q}(\gamma)\cong\mathbb{Q}(\gamma_i)\subseteq\mathbb{R}$($1\leqslant i\leqslant r_1$)为\emph{实嵌入}.而后 $r_2$ 对共轭域为虚域,并且 $\mathbb{Q}(\gamma_{r_1+j})=\mathbb{Q}(\gamma_{r_1+r_2+j})\not\subset\mathbb{R}$($1\leqslant j\leqslant r_2$).称这 $r_2$ 对嵌入 $\tau_i$($r_1+1\leqslant i\leqslant n$)为\emph{复嵌入},并且称 $\tau_{r_1+j}$ 和 $\tau_{r_1+r_2+j}$ 是彼此共轭的嵌入,记为 $\overline{\tau}_{r_1+j}=\tau_{r_1+r_2+j}$($1\leqslant j\leqslant r_2$).

如果 $L$ 中又有另一个元素 $\gamma'$ 使得 $L=\mathbb{Q}(\gamma')$,则 $\gamma'$ 在 $\mathbb{Q}$ 上的极小多项式 $f'(x)$ 的次数也是 $n=[L:\mathbb{Q}]$,并且 $f'(x)$ 的 $n$ 个根也有 $r_1$ 个实根和 $r_2$ 对虚根.这是由于嵌入 $\tau_1,\cdots,\tau_n$ 是域 $L$ 本身的特性,与生成元 $\gamma$ 的选取方式无关.换句话说,参数 $r_1$ 和 $r_2$ 是数域 $L$ 的不变量,$r_1+2r_2=n$(当然,扩张次数 $[L:\mathbb{Q}]=n$ 也是 $L$ 的不变量).
\end{remark}

\begin{example}\label{ex:ch1-1-1}
每个二次(数)域均可惟一地表示成 $\mathbb{Q}(\sqrt{d})$,其中 $d$ 是无平方因子整数(习题 2).当 $d>0$ 时它是实域,叫作实二次域;而当 $d<0$ 时叫作虚二次域.$\mathbb{Q}(\sqrt{d})/\mathbb{Q}$ 是伽罗瓦扩张,由于 $\sqrt{d}$ 的极小多项式为 $f(x)=x^2-d$,它的两个根是 $\pm\sqrt{d}$.所以伽罗瓦群 $\Gal(\mathbb{Q}(\sqrt{d})/\mathbb{Q})$ 为 $\{I,\sigma\}$,其中 $I$ 是恒等自同构,而 $\sigma(\sqrt{d})=-\sqrt{d}$.由于域 $\mathbb{Q}(\sqrt{d})$ 中每个元素惟一地表示成 $a+b\sqrt{d}$,其中 $a,b\in\mathbb{Q}$,可知 $\sigma(a+b\sqrt{d})=a-b\sqrt{d}$.不难看出,对于实二次域,$r_1=2$,$r_2=0$;而对于虚二次域,$r_1=0$,$r_2=1$.
\end{example}

\begin{example}\label{ex:ch1-1-2}
分圆域 $\mathbb{Q}(\zeta_{p^n})$,其中 $\zeta_{p^n}=\mathrm{e}^{2\pi\mathrm{i}/p^n}$ 是 $p^n$ 次本原单位根,而 $p$ 为素数,$n\geqslant 1$.以下简记 $\zeta=\zeta_{p^n}$.易知 $\zeta$ 是多项式
\[
  f(x)=x^{(p-1)p^{n-1}}+x^{(p-2)p^{n-1}}+\cdots+x^{p^{n-1}}+1
  =\frac{x^{p^n}-1}{x^{p^{n-1}}-1}\in\mathbb{Z}[x]
\]
的根.我们现在证明 $f(x)$ 是 $\mathbb{Q}[x]$ 中的不可约多项式.为此令
\[
  g(x)=f(x+1)=x^{(p-1)p^{n-1}}+c_{p^n-p^{n-1}-1}x^{p^n-p^{n-1}-1}+\cdots+c_1x+c_0\in\mathbb{Z}[x].
\]
由于
\[
  g(x)=\frac{(x+1)^{p^n}-1}{(x+1)^{p^{n-1}}-1}\equiv\frac{x^{p^n}}{x^{p^{n-1}}}=x^{p^n-p^{n-1}}\ (\mathrm{mod}\ p),
\]
从而 $p\mid c_i$($0\leqslant i\leqslant p^n-p^{n-1}-1$).进而 $c_0=g(0)=f(1)=p$,于是 $p^2\nmid c_0$.所以由 Eisenstein 判别准则(附录 A, (7))可知 $g(x)$ 是 $\mathbb{Q}[x]$ 中不可约多项式,从而 $f(x)$ 也是如此.这就表明 $f(x)$ 是 $\zeta$ 在 $\mathbb{Q}$ 上的极小多项式,并且 $[\mathbb{Q}(\zeta_{p^n}):\mathbb{Q}]=\deg f=p^n-p^{n-1}$.$\zeta$ 的全部共轭元素(即 $f(x)$ 的全部根)显然是
\[
  \zeta^i\quad (1\leqslant i\leqslant p^n,\ p\nmid i)
\]
(即为 $x^{p^n}-1$ 之根但不为 $x^{p^{n-1}}-1$ 之根者),它们均属于 $\mathbb{Q}(\zeta)$,从而 $\mathbb{Q}(\zeta)/\mathbb{Q}$ 是伽罗瓦扩张.令 $\sigma_i\in\Gal(\mathbb{Q}(\zeta)/\mathbb{Q})$,使得 $\sigma_i(\zeta)=\zeta^i$($1\leqslant i\leqslant p^n-1$,$p\nmid i$),则
\[
  \sigma_i\cdot\sigma_j(\zeta)=\sigma_i(\zeta^j)=\sigma_i(\zeta)^j=\zeta^{ij}=\sigma_{ij}(\zeta),
\]
这就表明 $\sigma_i\cdot\sigma_j=\sigma_{ij}$.作映射
\[
  \chi\colon\Gal(\mathbb{Q}(\zeta)/\mathbb{Q})\to(\mathbb{Z}/p^n\mathbb{Z})^\times,\qquad
  \chi(\sigma_i)=\overline{i},
\]
其中 $(\mathbb{Z}/p^n\mathbb{Z})^\times$ 表示环 $\mathbb{Z}/p^n\mathbb{Z}$ 的单位(乘法)群,而 $\overline{i}$ 表示剩余类 $i\ (\mathrm{mod}\ p^n)$.由上述不难看出 $\chi$ 是群的同构.但是当 $p\geqslant 3$ 时,从初等数论我们知道,乘法群 $(\mathbb{Z}/p^n\mathbb{Z})^\times$ 是由模 $p^n$ 的某个原根 $g$ 生成的 $p^n-p^{n-1}$ 阶循环群.从而伽罗瓦群 $\Gal(\mathbb{Q}(\zeta_{p^n})/\mathbb{Q})$ 当 $p\geqslant 3$ 时,是 $p^n-p^{n-1}$ 阶循环群,生成元为 $\sigma_g$.

当 $p^n\geqslant 3$ 时,$\zeta_{p^n}^i$($1\leqslant i\leqslant p^n$,$p\nmid i$)均不为实数.从而对于分圆域 $\mathbb{Q}(\zeta_{p^n})$($p^n\geqslant 3$)我们有 $r_1=0$,$r_2=\frac{1}{2}(p^n-p^{n-1})$.
\end{example}

\begin{example}\label{ex:ch1-1-3}
纯三次域 $\mathbb{Q}(\sqrt[3]{2})$.元素 $\sqrt[3]{2}$ 在 $\mathbb{Q}$ 上的极小多项式为 $x^3-2$.于是 $\sqrt[3]{2}$ 的共轭元素为 $\sqrt[3]{2}$,$\sqrt[3]{2}\omega$ 和 $\sqrt[3]{2}\omega^2$,其中 $\omega=\zeta_3$.由此可见 $\mathbb{Q}(\sqrt[3]{2})/\mathbb{Q}$ 不是伽罗瓦扩张,其正规闭包为
\[
  M=\mathbb{Q}(\sqrt[3]{2},\sqrt[3]{2}\omega,\sqrt[3]{2}\omega^2)=\mathbb{Q}(\sqrt[3]{2},\omega).
\]
不难算出 $[M:\mathbb{Q}]=6$ 而 $[\mathbb{Q}(\sqrt[3]{2}):\mathbb{Q}]=3$.并且对于域 $\mathbb{Q}(\sqrt[3]{2})$ 有 $r_1=r_2=1$.
\end{example}

\subsection{范与迹}\label{sec:ch1-s1-3}

设 $L/K$ 为数域扩张,$[L:K]=n$,$\sigma_i\colon L\to\mathbb{C}$($1\leqslant i\leqslant n$)是 $L$ 的 $n$ 个 $K$-嵌入.对于 $\alpha\in L$ 定义
\[
  \Norm_{L/K}(\alpha)=\prod_{i=1}^{n}\sigma_i(\alpha),\qquad
  \Tr_{L/K}(\alpha)=\sum_{i=1}^{n}\sigma_i(\alpha),
\]
分别称作是元素 $\alpha\in L$ 对于扩张 $L/K$ 的\emph{范}和\emph{迹}.从定义可知 $\Norm_{L/K}$ 和 $\Tr_{L/K}$ 有如下简单性质:
\begin{enumerate}[(a)]
  \item 对于 $\alpha,\beta\in L$,$\Norm_{L/K}(\alpha\beta)=\Norm_{L/K}(\alpha)\Norm_{L/K}(\beta)$,$\Tr_{L/K}(\alpha+\beta)=\Tr_{L/K}(\alpha)+\Tr_{L/K}(\beta)$;
  \item 对于 $\alpha\in K$,$\Norm_{L/K}(\alpha)=\alpha^n$,$\Tr_{L/K}(\alpha)=n\alpha$,其中 $n=[L:K]$.
\end{enumerate}

下面的定理 \ref{thm:1.3} 给出范和迹的一种计算方法.

\begin{theorem}\label{thm:1.3}
设 $L/K$ 为数域扩张,$[L:K]=n$,$\alpha\in L$,
\[
  f(x)=x^m-c_1x^{m-1}+\cdots+(-1)^mc_m\in K[x]
\]
是 $\alpha$ 在 $K$ 上的极小多项式,$m=[K(\alpha):K]$,则
\[
  \Norm_{L/K}(\alpha)=c_m^{n/m},\qquad
  \Tr_{L/K}(\alpha)=\frac{n}{m}c_1.
\]
\end{theorem}

\begin{proof}
设 $\alpha_1=\alpha,\alpha_2,\cdots,\alpha_m$ 是 $f(x)$ 的 $m$ 个根.由定义知
\[
  \Tr_{K(\alpha)/K}(\alpha)=\alpha_1+\alpha_2+\cdots+\alpha_m=c_1,
\]
\[
  \Norm_{K(\alpha)/K}(\alpha)=\alpha_1\alpha_2\cdots\alpha_m=c_m.
\]
根据定理 \ref{thm:1.2},每个 $K$-嵌入 $\tau_i\colon K(\alpha)\to\mathbb{C}$($\tau_i(\alpha)=\alpha_i$)均可扩充成 $[L:K(\alpha)]=n/m$ 个嵌入 $\sigma_{ij}\colon L\to\mathbb{C}$($1\leqslant j\leqslant n/m$).不难看出 $\{\sigma_{ij}\mid 1\leqslant i\leqslant m,\ 1\leqslant j\leqslant n/m\}$ 是彼此不同的,从而构成 $L$ 到 $\mathbb{C}$ 的全部 $K$-嵌入,于是
\begin{align*}
  \Norm_{L/K}(\alpha)
  &=\prod_{i=1}^{m}\prod_{j=1}^{n/m}\sigma_{ij}(\alpha)
   =\prod_{i=1}^{m}\prod_{j=1}^{n/m}\tau_i(\alpha)
   =\prod_{i=1}^{m}\prod_{j=1}^{n/m}\alpha_i\\
  &=(\alpha_1\cdots\alpha_m)^{n/m}=c_m^{n/m},
\end{align*}
\begin{align*}
  \Tr_{L/K}(\alpha)
  &=\sum_{i=1}^{m}\sum_{j=1}^{n/m}\sigma_{ij}(\alpha)
   =\sum_{i=1}^{m}\sum_{j=1}^{n/m}\tau_i(\alpha)
   =\sum_{i=1}^{m}\sum_{j=1}^{n/m}\alpha_i\\
  &=\frac{n}{m}(\alpha_1+\cdots+\alpha_m)=\frac{n}{m}c_1.
\end{align*}
\end{proof}

\begin{remark}
从定理 \ref{thm:1.3} 特别得到,对于每个 $\alpha\in L$,$\Norm_{L/K}(\alpha)$ 和 $\Tr_{L/K}(\alpha)$ 都是 $K$ 中的元素.再由性质 (a) 可知 $\Tr_{L/K}\colon L\to K$ 是加法群同态,而 $\Norm_{L/K}\colon L^\times\to K^\times$ 是乘法群同态,这里 $L^\times=L-\{0\}$.
\end{remark}

\begin{theorem}[\emph{传递公式}]\label{thm:1.4}
设 $L/M$,$M/K$ 均是数域扩张,$\alpha\in L$,则
\[
  \Norm_{L/K}(\alpha)=\Norm_{M/K}(\Norm_{L/M}(\alpha)),\qquad
  \Tr_{L/K}(\alpha)=\Tr_{M/K}(\Tr_{L/M}(\alpha)).
\]
\end{theorem}

\begin{proof}
令 $n=[L:M]$,$m=[M:K]$,$\sigma_1,\cdots,\sigma_n$ 为 $L$ 到 $\mathbb{C}$ 中 $n$ 个不同的 $M$-嵌入,$\tau_1,\cdots,\tau_m$ 为 $M$ 到 $\mathbb{C}$ 中 $m$ 个不同的 $K$-嵌入.取 $S$ 为扩张 $L/K$ 的正规闭包.令 $\widetilde{\sigma}_i$,$\widetilde{\tau}_j$ 分别为 $\sigma_i$ 和 $\tau_j$ 到 $S$ 上的一个扩充(定理 \ref{thm:1.2}),它们都是伽罗瓦群 $\Gal(S/K)$ 中的元素.从而 $(\widetilde{\tau}_j\widetilde{\sigma}_i)|_L$($1\leqslant i\leqslant n$,$1\leqslant j\leqslant m$)均是 $L$ 到 $\mathbb{C}$ 中的 $K$-嵌入.我们现在证明这 $nm$ 个 $K$-嵌入彼此不同:如果 $j\neq j'$,则存在 $b\in M$,使得 $\tau_j(b)\neq\tau_{j'}(b)$.于是
\[
  (\widetilde{\tau}_j\widetilde{\sigma}_i)|_L(b)=(\widetilde{\tau}_j\widetilde{\sigma}_i)(b)=\widetilde{\tau}_j(b)=\tau_j(b)\neq\tau_{j'}(b)=(\widetilde{\tau}_{j'}\widetilde{\sigma}_{i'})|_L(b).
\]
这就表明当 $j\neq j'$ 时 $(\widetilde{\tau}_j\widetilde{\sigma}_i)|_L\neq(\widetilde{\tau}_{j'}\widetilde{\sigma}_{i'})|_L$.如果 $j=j'$ 但是 $i\neq i'$,则存在 $c\in L$,使得 $\sigma_i(c)\neq\sigma_{i'}(c)$.于是
\[
  (\widetilde{\tau}_j\widetilde{\sigma}_i)|_L(c)=\widetilde{\tau}_j(\sigma_i(c))\neq\widetilde{\tau}_j(\sigma_{i'}(c))=(\widetilde{\tau}_j\widetilde{\sigma}_{i'})|_L(c),
\]
从而当 $i\neq i'$ 时 $(\widetilde{\tau}_j\widetilde{\sigma}_i)|_L\neq(\widetilde{\tau}_j\widetilde{\sigma}_{i'})|_L$.因此 $\{(\widetilde{\tau}_j\widetilde{\sigma}_i)|_L\mid 1\leqslant i\leqslant n,\ 1\leqslant j\leqslant m\}$ 就是 $L$ 到 $\mathbb{C}$ 中的全部 $nm=[L:K]$ 个 $K$-嵌入.从而对于每个 $\alpha\in L$,
\begin{align*}
  \Norm_{L/K}(\alpha)
  &=\prod_{1\leqslant j\leqslant m}\prod_{1\leqslant i\leqslant n}(\widetilde{\tau}_j\widetilde{\sigma}_i)|_L(\alpha)
   =\prod_{j=1}^{m}\widetilde{\tau}_j\Big(\prod_{i=1}^{n}\widetilde{\sigma}_i(\alpha)\Big)\\
  &=\prod_{j=1}^{m}\widetilde{\tau}_j(\Norm_{L/M}(\alpha))
   =\Norm_{M/K}(\Norm_{L/M}(\alpha)).
\end{align*}
对于迹可以类似地证明.
\end{proof}

\subsection{元素的判别式}\label{sec:ch1-s1-4}

设 $L/K$ 是数域的 $n$ 次扩张,$\sigma_1,\cdots,\sigma_n$ 是 $L$ 到 $\mathbb{C}$ 的 $n$ 个 $K$-嵌入,$\alpha_1,\cdots,\alpha_n$ 为 $L$ 中任意 $n$ 个元素.定义
\[
  d_{L/K}(\alpha_1,\cdots,\alpha_n)
  =\left|(\sigma_i(\alpha_j))_{\substack{1\leqslant i\leqslant n\\ 1\leqslant j\leqslant n}}\right|^2
\]
($|\,|$ 表示方阵的行列式),称作是元素 $\{\alpha_1,\cdots,\alpha_n\}$ 对于扩张 $L/K$ 的\emph{判别式}(从定义不难看出,这个判别式与元素 $\alpha_1,\cdots,\alpha_n$ 的次序是无关的).下面引理表明它是 $K$ 中的元素.

\begin{lemma}\label{lem:ch1-1}
$d_{L/K}(\alpha_1,\cdots,\alpha_n)=|(\Tr_{L/K}(\alpha_i\alpha_j))|$.
\end{lemma}

\begin{proof}
左边
\[
  =\begin{vmatrix}
    \sigma_1(\alpha_1)\cdots\sigma_n(\alpha_1)\\
    \vdots\\
    \sigma_1(\alpha_n)\cdots\sigma_n(\alpha_n)
  \end{vmatrix}
  \cdot
  \begin{vmatrix}
    \sigma_1(\alpha_1)\cdots\sigma_1(\alpha_n)\\
    \vdots\\
    \sigma_n(\alpha_1)\cdots\sigma_n(\alpha_n)
  \end{vmatrix}
  =\left|\Big(\sum_{k=1}^{n}\sigma_k(\alpha_i\alpha_j)\Big)\right|
  =\text{右边}.
\]
\end{proof}

判别式的第一个应用是它可用来判别 $L$ 中 $n$ 个元素 $\alpha_1,\cdots,\alpha_n$ 是否为向量空间 $L$ 的一组 $K$-基.

\begin{lemma}\label{lem:ch1-2}
$d_{L/K}(\alpha_1,\cdots,\alpha_n)\neq 0\iff\alpha_1,\cdots,\alpha_n$ 是 $K$-线性无关的.
\end{lemma}

\begin{proof}
$\Rightarrow$:如果 $\alpha_1,\cdots,\alpha_n$ 是 $K$-线性相关的,即存在不全为零的 $k_1,\cdots,k_n\in K$,使得 $k_1\alpha_1+\cdots+k_n\alpha_n=0$.于是
\[
  k_1\sigma_i(\alpha_1)+\cdots+k_n\sigma_i(\alpha_n)=\sigma_i(k_1\alpha_1+\cdots+k_n\alpha_n)=0\quad (1\leqslant i\leqslant n).
\]
这表明 $n$ 阶方阵 $(\sigma_i(\alpha_j))$ 的诸列是 $K$-线性相关的.从而 $d_{L/K}(\alpha_1,\cdots,\alpha_n)=|(\sigma_i(\alpha_j))|^2=0$.

$\Leftarrow$:如果 $d_{L/K}(\alpha_1,\cdots,\alpha_n)=0$,则 $|(\Tr_{L/K}(\alpha_i\alpha_j))|=0$(引理 \ref{lem:ch1-1}),从而方阵 $(\Tr_{L/K}(\alpha_i\alpha_j))$ 的 $n$ 行 $R_1,\cdots,R_n$ 是 $K$-线性相关的,其中
\[
  R_i=(\Tr_{L/K}(\alpha_i\alpha_1),\cdots,\Tr_{L/K}(\alpha_i\alpha_n))\quad (1\leqslant i\leqslant n).
\]
于是有不全为零的 $k_1,\cdots,k_n\in K$,使得 $k_1R_1+\cdots+k_nR_n=0$.如果 $\alpha_1,\cdots,\alpha_n$ 是 $K$-线性无关的,则 $\alpha=k_1\alpha_1+\cdots+k_n\alpha_n\neq 0$,而 $\Tr_{L/K}(\alpha\alpha_j)=0$($1\leqslant j\leqslant n$)因为它是向量 $k_1R_1+\cdots+k_nR_n=0$ 中的第 $j$ 个元素.由于当 $\alpha_1,\cdots,\alpha_n$ $K$-线性无关时,它们形成向量空间 $L$ 的一组 $K$-基,因此对于任何 $\beta\in L$ 均有 $\Tr_{L/K}(\alpha\beta)=0$.特别取 $\beta=\alpha^{-1}$(注意 $\alpha\neq 0$),则我们有
\[
  0=\Tr_{L/K}(\alpha\alpha^{-1})=\Tr_{L/K}(1)=[L:K]\neq 0.
\]
这一矛盾表明 $\alpha_1,\cdots,\alpha_n$ 是 $K$-线性相关的.
\end{proof}

\begin{lemma}\label{lem:ch1-3}
设 $L=K(\alpha)$,$[L:K]=n$,$f(x)$ 是 $\alpha$ 在 $K$ 上的极小多项式,$\alpha_1=\alpha,\alpha_2,\cdots,\alpha_n$ 为 $f(x)$ 的 $n$ 个复根,则
\[
  d_{L/K}(1,\alpha,\cdots,\alpha^{n-1})=\prod_{1\leqslant r<s\leqslant n}(\alpha_r-\alpha_s)^2
  =(-1)^{n(n-1)/2}\Norm_{L/K}(f'(\alpha)).
\]
\end{lemma}

\begin{proof}
上式左边
\[
  =\left|(\sigma_i(\alpha^j))_{\substack{1\leqslant i\leqslant n\\ 0\leqslant j\leqslant n-1}}\right|^2
  =\left|(\alpha_i^j)_{\substack{1\leqslant i\leqslant n\\ 0\leqslant j\leqslant n-1}}\right|^2,
\]
而后一行列式是 Vandermonde 行列式.从而可知原式左边等于 $\prod_{1\leqslant r<s\leqslant n}(\alpha_r-\alpha_s)^2$.进而,$(\alpha_r-\alpha_s)^2=-(\alpha_r-\alpha_s)(\alpha_s-\alpha_r)$,而满足 $1\leqslant r<s\leqslant n$ 的 $(r,s)$ 共有 $\frac{1}{2}n(n-1)$ 对.于是
\[
  \prod_{1\leqslant r<s\leqslant n}(\alpha_r-\alpha_s)^2
  =(-1)^{n(n-1)/2}\prod_{1\leqslant r\neq s\leqslant n}(\alpha_r-\alpha_s)
  =(-1)^{n(n-1)/2}\prod_{r=1}^{n}\prod_{\substack{s=1\\ s\neq r}}^{n}(\alpha_r-\alpha_s).
\]
但是 $f(x)=(x-\alpha_1)(x-\alpha_2)\cdots(x-\alpha_n)$,从而
$f'(\alpha_r)=\prod_{\substack{s=1\\ s\neq r}}^{n}(\alpha_r-\alpha_s)$.于是
\[
  \prod_{1\leqslant r<s\leqslant n}(\alpha_r-\alpha_s)^2
  =(-1)^{n(n-1)/2}\prod_{r=1}^{n}f'(\alpha_r)
  =(-1)^{n(n-1)/2}\Norm_{L/K}(f'(\alpha)).
\]
\end{proof}

\begin{remark}
对于 $L$ 中任意元素 $\alpha$,令
\[
  d_{L/K}(\alpha)=d_{L/K}(1,\alpha,\cdots,\alpha^{n-1}),\quad n=[L:K],
\]
称作是 $L$ 中元素 $\alpha$ 对于扩张 $L/K$ 的判别式.从引理 \ref{lem:ch1-2} 和引理 \ref{lem:ch1-3} 不难看出:
\[
  d_{L/K}(\alpha)\neq 0\iff 1,\alpha,\cdots,\alpha^{n-1}\text{ 是 }K\text{-线性无关的}
  \iff L=K(\alpha).
\]
当 $K=\mathbb{Q}$ 时,我们把 $d_{L/\mathbb{Q}}(\alpha)$ 也简写作 $d_L(\alpha)$.
\end{remark}

\begin{example}\label{ex:ch1-1-4}
考虑分圆域 $L=\mathbb{Q}(\omega)$,$\omega=\mathrm{e}^{2\pi\mathrm{i}/p^n}$,$p$ 为奇素数,$n\geqslant 1$.令 $s=p^n-p^{n-1}$,我们已经知道 $L$ 是 $s$ 次域,$\omega$ 在 $\mathbb{Q}$ 上的极小多项式为
\[
  f(x)=\frac{x^{p^n}-1}{x^{p^{n-1}}-1}=x^{(p-1)p^{n-1}}+x^{(p-2)p^{n-1}}+\cdots+x^{p^{n-1}}+1.
\]
于是
\[
  (x^{p^{n-1}}-1)f(x)=x^{p^n}-1,
\]
\[
  (x^{p^{n-1}}-1)f'(x)+p^{n-1}x^{p^{n-1}-1}f(x)=p^nx^{p^n-1}.
\]
因此 $f'(\omega)=p^n/\big(\omega(\omega^{p^{n-1}}-1)\big)$.于是由引理 \ref{lem:ch1-3} 可知
\[
  d_L(\omega)=(-1)^{s(s-1)/2}\Norm(p^n)\Big/\Big(\Norm(\omega)\Norm(\omega^{p^{n-1}}-1)\Big),
\]
其中 $\Norm=\Norm_{L/\mathbb{Q}}$.易知 $\Norm(p^n)=p^{ns}$,
\[
  \Norm(\omega)=\prod_{\substack{k=1\\ p\nmid k}}^{p^{n-1}}\omega^k=\prod_{\substack{k=1\\ p\nmid k}}^{(p^{n-1})/2}(\omega^k\cdot\omega^{-k})=1.
\]
剩下只需计算 $\Norm(\omega^{p^{n-1}}-1)$.令 $\zeta=\omega^{p^{n-1}}=\mathrm{e}^{2\pi\mathrm{i}/p}$,则 $\zeta$ 在 $\mathbb{Q}$ 上的极小多项式为 $g(x)=x^{p-1}+x^{p-2}+\cdots+x+1$,从而 $\zeta-1$ 在 $\mathbb{Q}$ 上的极小多项式为 $g(x+1)=x^{p-1}+\cdots+p$.令 $M=\mathbb{Q}(\zeta)$,则 $[M:\mathbb{Q}]=p-1$,于是 $[L:M]=s/(p-1)=p^{n-1}$.从而
\[
  \Norm_{L/\mathbb{Q}}(\zeta-1)=\Norm_{M/\mathbb{Q}}(\Norm_{L/M}(\zeta-1))
  =\Norm_{M/\mathbb{Q}}(\zeta-1)^{p^{n-1}}
  =p^{p^{n-1}}.
\]
从而最后得到
\[
  d_L(\omega)=(-1)^{s(s-1)/2}p^{ns}\big/p^{p^{n-1}}
  =(-1)^{s(s-1)/2}p^{p^{n-1}(np-n-1)}.
\]
\end{example}

请读者验证,当 $p=2$ 时这个公式仍然成立.

\subsection{单位根}\label{sec:ch1-s1-5}

数域 $K$ 中的乘法有限阶元素 $\omega\neq 0$ 叫作是 $K$ 中的\emph{单位根}.如果 $\omega^n=1$,则称 $\omega$ 为 $n$ 次单位根.如果 $\omega$ 的乘法阶数恰好是 $n$,即 $n$ 是满足 $\omega^n=1$ 的最小正整数,则称 $\omega$ 是 $n$ 次\emph{本原}单位根.数域 $K$ 中的单位根全体 $W_K$ 显然形成乘法群,称作是数域 $K$ 的\emph{单位根群}.我们现在要证明 $W_K$ 是有限循环群.首先证明:

\begin{lemma}\label{lem:ch1-4}
$W_K$ 是有限群.
\end{lemma}

\begin{proof}
设 $\omega$ 是数域 $K$ 中的 $n$ 次单位根.令
\[
  f(x)=x^m+c_1x^{m-1}+\cdots+c_m
\]
是 $\omega$ 在 $\mathbb{Q}$ 上的极小多项式,$f(x)$ 的全部根为 $\omega=\omega_1,\omega_2,\ldots,\omega_m$.由于 $\omega^n=1$,从而 $f(x)\mid x^n-1$,于是 $\omega$ 的每个共轭元素也均满足 $\omega_i^n=1$,即 $\omega_i$ 均是 $n$ 次单位根,于是 $|\omega_i|=1$($1\leqslant i\leqslant m$).由韦达定理可知 $|c_i|\leqslant\binom{m}{i}$($1\leqslant i\leqslant m$).另一方面,$f(x)$ 是 $x^n-1$ 的首 $1$(即最高次项系数为 $1$)的多项式因子,由此不难证明 $f(x)\in\mathbb{Z}[x]$,即系数 $c_i$ 均为有理整数.进而,由于 $\omega\in K$,可知
\[
  m=\deg f=[\mathbb{Q}(\omega):\mathbb{Q}]\leqslant[K:\mathbb{Q}].
\]
但是次数 $\leqslant[K:\mathbb{Q}]$ 而有理整系数又有界 $|c_i|\leqslant\binom{m}{i}\leqslant\binom{[K:\mathbb{Q}]}{i}$($1\leqslant i\leqslant m$)的多项式 $f(x)$ 只有有限多个,从而它们的根也只有有限多个.这就表明每个数域 $K$ 中均只有有限多个单位根,即 $W_K$ 是有限群.
\end{proof}

\begin{lemma}\label{lem:ch1-5}
任意域中的有限乘法群均是循环群.
\end{lemma}

\begin{proof}
设 $W$ 是域 $K$ 中的有限乘法群.令 $|W|=n$.当 $n=1$ 时引理显然正确.若 $n\geqslant 2$,令 $n=p_1^{\alpha_1}p_2^{\alpha_2}\cdots p_s^{\alpha_s}$,其中 $p_1,\cdots,p_s$ 是不同的素数,$\alpha_i\geqslant 1$.由于 $W$ 是乘法 Abel 群,根据有限 Abel 群的结构定理(附录 A, (1)),
\[
  W=W_1\times W_2\times\cdots\times W_s\ (\text{直积}),
\]
其中 $W_i$ 是 $W$ 的 $p_i^{\alpha_i}$ 阶 Sylow 子群($1\leqslant i\leqslant s$).于是 $W_i$ 中每个元素 $\alpha$ 均满足 $\alpha^{p_i^{\alpha_i}}=1$.但是在域 $K$ 中多项式 $x^{p_i^{\alpha_i}-1}-1$ 至多有 $p_i^{\alpha_i}-1$ 个根,从而在 $W_i$ 中满足 $x^{p_i^{\alpha_i}-1}=1$ 的元素 $x$ 也至多有 $p_i^{\alpha_i}-1$ 个.这就表明 $W_i$ 中存在 $p_i^{\alpha_i}$ 阶元素 $w_i$($1\leqslant i\leqslant s$).由于 $p_1^{\alpha_1},\cdots,p_s^{\alpha_s}$ 两两互素,因此 $W$ 中元素 $w=w_1w_2\cdots w_s$ 的阶为 $p_1^{\alpha_1}\cdots p_s^{\alpha_s}=n=|W|$.这就表明 $W$ 是循环群.
\end{proof}

从以上两个引理立刻得出:

\begin{theorem}\label{thm:1.5}
数域 $K$ 中的单位根群 $W_K$ 是有限循环群.
\end{theorem}

通常以 $w_K$ 表示 $K$ 中单位群 $W_K$ 的阶.

\begin{example}\label{ex:ch1-1-5}
如果 $K$ 为实域,则 $W_K=\{\pm1\}$,$w_K=2$.
\end{example}

\begin{example}\label{ex:ch1-1-6}
\begin{sloppypar}
设 $K=\mathbb{Q}(\sqrt{-d})$ 为虚二次域,$d$ 为无平方因子的正整数,则当 $d=1$ 时,$W_K=\{\pm1,\pm\sqrt{-1}\}$,$w_K=4$;当 $d=3$ 时,$W_K=\{\pm1,\pm\omega,\pm\omega^2\}$,$\omega=\mathrm{e}^{2\pi\mathrm{i}/3}=\frac{1}{2}(-1+\sqrt{-3})$,$w_K=6$;而当 $d>3$ 时,$W_K=\{\pm1\}$,$w_K=2$.这是因为:如果 $\zeta_n\in K$ 并且 $n$ 有因子 $p^m$,则
\end{sloppypar}
\[
  K\supseteq\mathbb{Q}(\zeta_n)\supseteq\mathbb{Q}(\zeta_{p^m}),
\]
从而
\[
  2=[K:\mathbb{Q}]\geqslant[\mathbb{Q}(\zeta_n):\mathbb{Q}]\geqslant[\mathbb{Q}(\zeta_{p^m}):\mathbb{Q}]=p^{m-1}(p-1).
\]
从而当 $p^{m-1}(p-1)\geqslant 3$ 时,$K$ 中不可能包含 $\zeta_n$,于是也就不能包含任何 $n$ 次本原单位根.所以能够在虚二次域 $K$ 中的只能是 $6$ 次和 $3$ 次本原单位根(即上面的 $\pm\omega$ 和 $\pm\omega^2$,此时 $K=\mathbb{Q}(\sqrt{-3})$),$4$ 次本原单位根(即 $\pm\sqrt{-1}$,此时 $K=\mathbb{Q}(\sqrt{-1})$)和 $\pm1$.
\end{example}

对于每个正整数 $n$,$\zeta_n=\mathrm{e}^{2\pi\mathrm{i}/n}$ 是 $n$ 次本原单位根,而 $\zeta_n^k$($1\leqslant k\leqslant n$,$(k,n)=1$)就是全部 $n$ 次本原单位根.它们共有 $\varphi(n)$ 个,其中 $\varphi(n)$ 是\emph{欧拉函数},表示 $1,2,\cdots,n$ 当中与 $n$ 互素的数的个数.$\varphi(n)$ 是积性函数,即当 $(n,m)=1$ 时,$\varphi(nm)=\varphi(n)\varphi(m)$.由此及 $\varphi(p^n)=p^n-p^{n-1}=p^n(1-\frac{1}{p})$ 就得到公式
\[
  \varphi(n)=n\prod_{p\mid n}\Big(1-\frac{1}{p}\Big).
\]

现在我们考虑一般的分圆域 $\mathbb{Q}(\zeta_n)$,$\zeta_n=\mathrm{e}^{2\pi\mathrm{i}/n}$,$n$ 是任意正整数.为了计算这个域的次数,我们需要决定 $\zeta_n$ 的全部共轭元素.

\begin{lemma}\label{lem:ch1-6}
与 $\zeta_n$ 共轭的全部元素即是 $\varphi(n)$ 个 $n$ 次本原单位根.
\end{lemma}

\begin{proof}
$\zeta_n$ 的每个共轭元素均是某个域嵌入之下元素 $\zeta_n$ 的象,由此不难看出 $\zeta_n$ 的每个共轭元素均是 $n$ 次本原单位根.问题在于我们还需证明每个 $n$ 次本原单位根均与 $\zeta_n$ 共轭.为此我们只需证明:对于每个 $n$ 次本原单位根 $\omega$ 和素数 $p$,如果 $p\nmid n$,则 $\omega^p$ 必然与 $\omega$ 共轭,因为如果这件事成立,就可以对于 $k$ 归纳证明每个 $n$ 次本原单位根 $\zeta_n^k$($1\leqslant k\leqslant n$,$(n,k)=1$)均与 $\zeta_n$ 共轭.

设 $f(x)$ 为 $\omega$ 在 $\mathbb{Q}$ 上的极小多项式.由于 $\omega^n=1$,从而 $f(x)\mid x^n-1$.令 $x^n-1=f(x)g(x)$.因为 $f$ 和 $g$ 均是 $x^n-1$ 的首 $1$ 多项式因子,因此它们事实上均属于 $\mathbb{Z}[x]$(习题 10).现在假定 $\omega^p$ 与 $\omega$ 不共轭,则 $f(\omega^p)\neq 0$,于是 $g(\omega^p)=0$(因为 $(\omega^p)^n-1=0$).从而 $\omega$ 是多项式 $g(x^p)$ 的根.于是 $f(x)\mid g(x^p)$.令 $g(x^p)=f(x)h(x)$.由于 $g(x^p),f(x)\in\mathbb{Z}[x]$,从而首 $1$ 多项式 $h(x)$ 也属于 $\mathbb{Z}[x]$.在环的自然同态 $\mathbb{Z}[x]\to(\mathbb{Z}/p\mathbb{Z})[x]$ 之下,
\[
  \overline{f}(x)\mid\overline{g}(x^p)=\overline{g}(x)^p.
\]
\begin{sloppypar}
由于 $(\mathbb{Z}/p\mathbb{Z})[x]$ 为惟一因子分解整环,从而上式表明 $(\overline{f}(x),\overline{g}(x))\neq 1$.因此存在多项式 $\overline{k}(x)\in(\mathbb{Z}/p\mathbb{Z})[x]$,$\deg\overline{k}(x)\geqslant 1$,使得 $\overline{k}(x)\mid(\overline{f}(x),\overline{g}(x))$.于是 $\overline{k}^2(x)\mid\overline{f}\cdot\overline{g}=x^n-\overline{1}$.但这是不可能的,因为当 $p\nmid n$ 时 $(\mathbb{Z}/p\mathbb{Z})[x]$ 中多项式 $x^n-\overline{1}$ 没有重根(附录 A, (10)).这一矛盾表明 $\omega^p$ 与 $\omega$ 共轭.从而也就证明了引理 \ref{lem:ch1-6}.
\end{sloppypar}
\end{proof}

\begin{theorem}\label{thm:1.6}
\begin{enumerate}[(1)]
  \item $[\mathbb{Q}(\zeta_n):\mathbb{Q}]=\varphi(n)$,并且 $\zeta_n=\mathrm{e}^{2\pi\mathrm{i}/n}$ 在 $\mathbb{Q}$ 上的极小多项式为
        \[
          \Phi_n(x)=\prod_{\substack{1\leqslant k\leqslant n\\ (n,k)=1}}(x-\zeta_n^k)\in\mathbb{Z}[x]
        \]
        ($\Phi_n(x)$ 叫作\emph{分圆多项式}).
  \item $\mathbb{Q}(\zeta_n)/\mathbb{Q}$ 为伽罗瓦扩张,并且其伽罗瓦群自然同构于 $(\mathbb{Z}/n\mathbb{Z})^\times$.
\end{enumerate}
\end{theorem}

\begin{proof}
(1) 由引理 \ref{lem:ch1-6} 推出.

(2) 由于 $\zeta_n$ 的全部共轭元素均属于 $K=\mathbb{Q}(\zeta_n)$,从而 $K/\mathbb{Q}$ 是伽罗瓦扩张.令 $\sigma_k\in\Gal(K/\mathbb{Q})$,使得 $\sigma_k(\zeta_n)=\zeta_n^k$,$(n,k)=1$,则
\[
  \sigma_k\sigma_{k'}(\zeta_n)=\sigma_k(\zeta_n^{k'})=\zeta_n^{k'k}=\sigma_{kk'}(\zeta_n),
\]
因此 $\sigma_k\sigma_{k'}=\sigma_{kk'}$.所以映射
\[
  \chi\colon\Gal(K/\mathbb{Q})\to(\mathbb{Z}/n\mathbb{Z})^\times,\qquad
  \chi(\sigma_k)=\overline{k}
\]
是群的满同态.但是由引理 \ref{lem:ch1-6},
\[
  |\Gal(K/\mathbb{Q})|=[K:\mathbb{Q}]=\varphi(n)=|(\mathbb{Z}/n\mathbb{Z})^\times|,
\]
从而 $\chi$ 是同构.
\end{proof}

当 $n\equiv 2\ (\mathrm{mod}\ 4)$ 时,$\zeta_n=-\zeta_n^{(2+n)/2}=-\zeta_{n/2}^{(n+2)/4}$.因此 $\mathbb{Q}(\zeta_n)=\mathbb{Q}(\zeta_{n/2})$.所以通常对于分圆域 $\mathbb{Q}(\zeta_n)$,可规定 $n\not\equiv 2\ (\mathrm{mod}\ 4)$.在这一规定下,可以证明当 $n\neq n'$ 时 $\mathbb{Q}(\zeta_n)\neq\mathbb{Q}(\zeta_{n'})$(习题 8),并且也不难决定分圆域的单位根群(习题 9).

% ============ 习题 ============
\begin{Exercises}
  \item 证明:代数数集合是可数的,而超越数集合是不可数的.
  \item (a) 求证每个二次(数)域均可表达成 $\mathbb{Q}(\sqrt{d})$,其中 $d$ 是无平方因子整数;
        (b) 如果 $d$ 和 $d'$ 均是无平方因子整数,并且 $d\neq d'$,则 $\mathbb{Q}(\sqrt{d})\neq\mathbb{Q}(\sqrt{d'})$;
        (c) 二次域 $K$ 必然是 $\mathbb{Q}$ 的伽罗瓦扩张,试求其伽罗瓦群.
  \item (a) 求下列代数数的次数和(在 $\mathbb{Q}$ 上的)极小多项式:
        $\sqrt{2}+\sqrt{3}+\sqrt{5}$,$\sqrt{2+\sqrt{2}}$,$\sqrt[3]{2}+\omega$(其中 $\omega=\mathrm{e}^{2\pi\mathrm{i}/3}$);
        (b) 求元素 $\sqrt{2}+\sqrt{3}+\sqrt{5}$ 对 $\mathbb{Q}(\sqrt{30})$ 的全部共轭元素.
  \item 证明代数数全体组成的集合 $\Omega$ 是域,并且是有理数域 $\mathbb{Q}$ 的无限(次)代数扩域.
  \item 对于数域 $L=\mathbb{Q}(\sqrt{1+\sqrt{2}},\sqrt{1-\sqrt{2}})$ 和 $L=\mathbb{Q}(\sqrt{2}+\sqrt{3},\sqrt{2}-\sqrt{5})$,求元素 $\gamma$,使得 $L=\mathbb{Q}(\gamma)$.
  \item 设 $f(x)\in K[x]$ 是数域 $K$ 上的 $n$ 次不可约首 $1$ 多项式,$\alpha_1,\cdots,\alpha_n$ 是它的 $n$ 个根,称
        \[
          d(f)=\prod_{1\leqslant r<s\leqslant n}(\alpha_r-\alpha_s)^2
        \]
        是多项式 $f(x)$ 的判别式.
        (a) 求证 $d(f)$ 是 $K$ 中元素;
        (b) 设 $f(x)=x^n+a$,$a\in\mathbb{Q}$,$\sqrt[n]{-a}\notin\mathbb{Q}$,求证 $d(f)=(-1)^{n(n-1)/2}n^n\alpha^{n-1}$;
        (c) 设 $f(x)=x^n+ax+b$ 是 $\mathbb{Q}[x]$ 中不可约多项式,求证
        \[
          d(f)=(-1)^{n(n-1)/2}\big[(-1)^{n-1}(n-1)^{n-1}a^n+n^nb^{n-1}\big]
        \]
        (注:当 $n=2$ 和 $3$ 时,$d(f)$ 分别为 $a^2-4b$ 和 $-(4a^3+27b^2)$,这就是 $2$ 次和 $3$ 次多项式通常所谓的判别式).
  \item 求证:一个代数整数是单位根的充要条件是它的每个共轭元素的绝对值均是 $1$.
  \item 设 $n$ 和 $m$ 为正整数,$\mathbb{Q}(\zeta_n)=Q_n$,
        (1) $Q_n\cap Q_m=Q_{(n,m)}$,$Q_nQ_m=Q_{[n,m]}$;
        (2) 如果 $n\neq n'$,$n\not\equiv 2\ (\mathrm{mod}\ 4)$,$n'\not\equiv 2\ (\mathrm{mod}\ 4)$,求证 $\mathbb{Q}(\zeta_n)\neq\mathbb{Q}(\zeta_{n'})$.
  \item 令 $K=\mathbb{Q}(\zeta_n)$,则
        (a) 当 $n\equiv 1\ (\mathrm{mod}\ 2)$ 时,$W_K=\{\zeta_{2n}^k\mid 0\leqslant k\leqslant 2n-1\}$;
        (b) 当 $n\equiv 0\ (\mathrm{mod}\ 4)$ 时,$W_K=\{\zeta_n^k\mid 0\leqslant k\leqslant n-1\}$.
  \item 如果 $f(x)$ 是 $\mathbb{Z}[x]$ 中的首 $1$ 多项式,而 $g(x)\in\mathbb{Q}[x]$ 是 $f(x)$ 的首 $1$ 多项式因子,求证 $g(x)\in\mathbb{Z}[x]$.
  \item 设 $f(x)$ 为 $\mathbb{Z}[x]$ 中不可约首 $1$ 多项式,$\alpha_1,\cdots,\alpha_n$ 是它的 $n$ 个根,令 $s_m=\sum_{i=1}^{n}\alpha_i^m$,求证
        \[
          d(f)=\begin{vmatrix}
            s_0 & s_1 & \cdots & s_{n-1}\\
            s_1 & s_2 & \cdots & s_n\\
            \vdots & \vdots & & \vdots\\
            s_{n-1} & s_n & \cdots & s_{2n-2}
          \end{vmatrix}.
        \]
  \item 设 $f(x)$ 是 $\mathbb{Q}[x]$ 中不可约三次首 $1$ 多项式,求证:如果 $d(f)>0$,则 $f(x)$ 有 $3$ 个实根;如果 $d(f)<0$,则 $f(x)$ 只有 $1$ 个实根.
  \item 设 $L/K$ 是数域的扩张.对于 $\alpha\in L$,定义映射
        \[
          \varphi_\alpha\colon L\to L,\quad \varphi_\alpha(\beta)=\alpha\beta.
        \]
        (a) 求证 $\varphi_\alpha$ 是 $K$-向量空间 $L$ 中的线性变换;
        (b) 如果 $A_\alpha$ 是线性变换 $\varphi_\alpha$ 对于向量空间 $L$ 的任意一组 $K$-基的变换方阵,求证 $\Norm_{L/K}(\alpha)=|A_\alpha|$,$\Tr_{L/K}(\alpha)=\Tr(A_\alpha)$(其中 $\Tr(A)$ 表示方阵 $A$ 的迹).
\end{Exercises}
